\section{Introduction}
\label{sec:introduction}

An agent that shops on someone's behalf faces the same question after every merchant
it queries: buy this, or keep looking? Looking again might find a better price. It
also spends wall-clock time, model tokens, API calls, and money, and the offer in
hand may not survive the detour.

In practice this question is answered by a constant. Query three merchants. Spend at
most thirty seconds. Accept anything under a threshold. These rules are easy to ship
and impossible to defend, because nobody has measured when they are good enough.
Rather than argue the point from theory, we measure it.

\paragraph{What we do.}
We collect two dated snapshots of a large set of merchants that publish
machine-readable catalogs, turn them into replayable search episodes, and run eight
stopping policies against the same frozen panels. The policies range from accepting
the first available offer to an exact reservation-price rule, and include an offline
oracle and a deliberately broken control.

\paragraph{What we find.}
Three results, in increasing order of practical usefulness.

\begin{enumerate}
  \item \textbf{Catalog listings differ across the two dates.} Between the two snapshots,
    \ResultDelistingRate{} of tracked listings disappeared and
    \ResultPriceChangeRate{} of the survivors changed price. Repeating a scan minutes
    apart reproduces it exactly, so this is catalog change rather than noise in the
    instrument. It does not establish checkout availability, free recall, or the
    irrecoverability of an earlier offer. (Claim \texttt{OFFER-EPHEMERALITY-001}.)
  \item \textbf{Adaptive stopping wins, but barely.} On
    \ResultHeldOutEpisodeCount{} held-out episodes the adaptive rule beats the
    strongest tuned fixed rule, with the entire bootstrap interval below zero and
    \ResultHardBudgetViolations{} budget violations. The effect is a fraction of a
    percent of purchase cost. We report the small size as prominently as the
    significance. (Claim \texttt{STOPPING-ADVANTAGE-001}.)
  \item \textbf{The reason it is small is the useful part.} In this corpus the median
    product is carried by \ResultMedianMerchantCount{} merchants at a price ratio of
    \ResultMedianPriceDispersionRatio{}. When merchants charge the same price, no
    stopping rule can help. A sweep calibrated to the measured churn and dispersion
    shows adaptive stopping helping only once prices genuinely differ and the budget
    allows acting on what a query reveals. Fixed rules win in some grid cells at both
    constrained and full budgets, so the grid does not establish one shared
    mechanism. (Claim \texttt{NO-ADVANTAGE-REGION-001}.)
\end{enumerate}

\paragraph{What to take away.}
The deliverable is not merely the dynamic program; that is standard backward
induction, deferred to Appendix~\ref{app:proofs}. The deliverable is the decision
framework and Table~\ref{tab:decision}: given the price spread, merchant count, and
budget observed up front, it determines when adaptive stopping provides a genuine
economic advantage and when simpler search heuristics suffice.

\paragraph{Scope.}
Every number in this paper is regenerated from the collected data by the build; none
are typed by hand. Claims are registered with an executable evaluator before being
reported, and results derived from simulation are labelled as such wherever they
appear. The primary comparison uses commit-or-continue semantics: continuing does
not reserve the observed catalog offer. Recall or revalidation requires an additional
state and action model and is not evaluated here.

\paragraph{Ongoing collection.}
The two-snapshot analysis reported here is the current evidence release. We are
continuing daily re-observation of the same SKU--merchant panel and weekly discovery
of new multi-merchant episodes. Those observations will support a future analysis of
catalog-offer survival and repricing across minutes, hours, and days; they are not
used for any result reported in this version.
